\documentclass[11pt,a4paper]{article}

% Codificación y lenguaje
\usepackage[utf8]{inputenc}
\usepackage[T1]{fontenc}
\usepackage[spanish,es-nodecimaldot]{babel}

% Presentación
\usepackage[a4paper,margin=2.7cm]{geometry}
\usepackage{parskip}
\usepackage{booktabs}
\usepackage{array}
\usepackage{longtable}
\usepackage{xcolor}
\usepackage{fancyhdr}
\usepackage{lastpage}
\usepackage{amsmath}
\usepackage{listings}
\usepackage[hidelinks]{hyperref}

\definecolor{codebg}{RGB}{247,248,250}
\definecolor{codecomment}{RGB}{70,110,70}
\definecolor{codekeyword}{RGB}{0,70,140}
\definecolor{codestring}{RGB}{150,60,30}

\lstdefinelanguage{VerilogCustom}{
  morekeywords={module,endmodule,input,output,inout,wire,reg,parameter,localparam,
    always,begin,end,case,endcase,default,if,else,assign,signed,integer,for,
    initial,display,finish},
  sensitive=true,
  morecomment=[l]{//},
  morecomment=[s]{/*}{*/},
  morestring=[b]"
}

\lstset{
  language=VerilogCustom,
  basicstyle=\ttfamily\small,
  keywordstyle=\color{codekeyword}\bfseries,
  commentstyle=\color{codecomment}\itshape,
  stringstyle=\color{codestring},
  backgroundcolor=\color{codebg},
  frame=single,
  rulecolor=\color{black!20},
  numbers=left,
  numberstyle=\tiny\color{black!50},
  numbersep=8pt,
  breaklines=true,
  columns=fullflexible,
  keepspaces=true,
  showstringspaces=false,
  tabsize=4
}

\pagestyle{fancy}
\fancyhf{}
\lhead{Arquitectura de Computadoras}
\rhead{Trabajo Final -- Procesador RISC-V}
\cfoot{\thepage\ de \pageref{LastPage}}
\setlength{\headheight}{14pt}

\newcommand{\verilog}[1]{\texttt{\detokenize{#1}}}

\title{\textbf{Diseño e implementación de un procesador RISC-V segmentado}\\[0.4em]
       \large Trabajo Final}
\author{Javier Siñuka y David Trujillo}
\date{\today}

\begin{document}

\begin{titlepage}
  \hypersetup{pageanchor=false}
  \centering
  \vspace*{1.2cm}
  {\Large Universidad \par}
  \vspace{0.3cm}
  {\large Arquitectura de Computadoras \par}
  \vspace{2.2cm}
  {\LARGE\bfseries Diseño e implementación de un procesador RISC-V segmentado\par}
  \vspace{0.5cm}
  {\large Trabajo Final\par}
  \vfill
  \begin{tabular}{rl}
    Estudiantes: & Javier Siñuka y David Trujillo \\
    Fecha: & \today
  \end{tabular}
  \vfill
\end{titlepage}

\hypersetup{pageanchor=true}
\clearpage
\pagenumbering{roman}
\tableofcontents
\clearpage
\pagenumbering{arabic}

\begin{abstract}
% Resumir el objetivo, la arquitectura implementada y los resultados obtenidos.
\end{abstract}

\section{Introducción}

\section{Objetivos}

\section{Arquitectura general}

\section{Pipeline del procesador}

\subsection{Etapas IF, ID, EX, MEM y WB}

\subsection{Registros intermedios}

\section{Conjunto de instrucciones implementado}

\section{Riesgos y mecanismos de resolución}

\section{Memorias y carga dinámica de programas}

\section{Unidad de depuración e interfaz UART}

\section{Modos de operación}

\section{Interfaz de usuario}

\section{Implementación en FPGA y temporización}

\section{Simulación y verificación}

\section{Conclusiones}

\end{document}

